\section{Implementación y validación en la placa}\label{cap:implementacion}
\subsection{Trazabilidad del entorno de desarrollo}
El proyecto combina herramientas de entrenamiento, conversión, compilación y programación. Se conservará un manifiesto por ensayo. Las versiones citadas en distintas etapas no se fusionarán en un único entorno supuesto: los prototipos usaron configuraciones diferentes.

\begin{table}[H]\centering\small
\caption{Herramientas mencionadas en el desarrollo y verificación pendiente.}
\begin{tabularx}{\linewidth}{p{4cm}Y}\toprule
Componente & Registro de trabajo\\\midrule
ST Edge AI Core & 4.0.1; conservar salida de versión y opciones\\
STM32CubeIDE & 1.18.1 en los ensayos registrados\\
STM32Cube.AI & 12.0.1-RC2 en registros de desarrollo\\
TensorFlow / Keras / TFMOT & Exportar versiones efectivas de cada entorno; no asumir compatibilidad entre registros\\
Biblioteca de ejecución & Conservar archivo y huella de la variante CM55 utilizada\\
Sistema anfitrión y compilador & Registrar versión, arquitectura y opciones de optimización\\\bottomrule
\end{tabularx}\end{table}

\subsection{Exportación, análisis y generación}
El artefacto de entrada se valida antes de pasar a ST Edge AI. Se conservarán el reporte de análisis, código generado, descriptores de memoria y mensajes de conversión. La etapa se considera satisfactoria cuando el modelo compilado conserva la calidad necesaria, no únicamente cuando la compilación termina sin errores.

El análisis registrará qué operadores se ejecutan en hardware y cuáles en software. Si la herramienta promueve una red INT16x8 a Float32, se genera una fila distinta en la comparación de implementaciones. La latencia de esa fila pertenece a Float32 compilado a partir de INT16x8, y no a una inferencia INT16x8 nativa.

\subsection{Proyecto STM32 y arranque}
La integración utiliza el flujo de aplicación y arranque configurado para STM32N6. En los ensayos se trabajó con el flujo \termino{Trusted}, incluyendo aplicación y datos de red. Direcciones de memoria, cargador externo y archivos de programación se obtendrán del proyecto generado y de su mapa de memoria, sin fijarlas por analogía con otro ejemplo.

El anexo técnico conservará la función de entrada de inferencia, inicialización de búferes, tratamiento de errores y configuración UART. Se documentarán también la primera etapa de arranque y la relación entre los binarios programados. No se copiarán módulos entre proyectos de arranque sin verificar su función.

\subsection{Interfaces y contrato numérico}
Cada entrada tendrá nombre, unidad, orden, forma, tipo, escala y punto cero. En radiografías se fijará el orden NHWC y la normalización. En MPC se fijará explícitamente la permutación de estados, el tratamiento del ángulo y la conversión de salida a voltios.

La validación compara el resultado en placa con el del intérprete para la misma entrada serializada. Se registrará el error antes y después de descuantizar. Este control ayuda a separar un error de comunicación o de orden de canales de un efecto real de la cuantización.

\subsection{Instrumentación temporal}
El temporizador DWT o el mecanismo temporal disponible permite medir ciclos de ejecución. La conversión requiere una frecuencia efectiva conocida:
\begin{equation}t_{\mathrm{inf}}=\frac{n_{\mathrm{ciclos}}}{f_{\mathrm{CPU}}}.\end{equation}
En una inferencia NPU, el intervalo CPU debe abarcar la finalización de la operación asíncrona para representar el tiempo transcurrido completo. Se documentarán el desbordamiento del contador, el calentamiento y el costo del instrumento. Los mensajes de depuración se excluyen del intervalo cuando se pretende medir sólo la inferencia.

\subsection{UART, repetición e integridad de los registros}
Los ensayos utilizaron COM4 a 115200 baudios como configuración de trabajo del computador. El puerto no es parte de la definición científica del método y debe poder modificarse. Se conservarán número de secuencia, tiempos, estado de recepción y el registro íntegro de cada corrida.

Una inferencia de aproximadamente 0,34 ms, un RTT cercano a 5 ms y un período de 160 ms son tres magnitudes compatibles y distintas. La diferencia incluye transmisión, planificación del computador y lógica de aplicación. Para demostrar operación en tiempo real del lazo completo se medirá el ciclo efectivo y el número de incumplimientos.
\pendiente{Adjuntar proyecto final, mapa de memoria, exportación de perfil, huellas de los modelos y ejemplos mínimos de llamada a inferencia. El código de producción debe provenir del proyecto validado.}
